% \subsection{CommandRouter} 
% \begin{algorithm}[H]
% \caption{CommandRouter (\emph{Interface})}
% \begin{algorithmic}[1]
% \Procedure{Route}{command : Command\textlangle REQ\textrangle} 
%     \State \textbf{Input:} command of type \textlangle REQ\textrangle
%     \State \textbf{Output:} CommandHandler of type\textlangle REQ, RES\textrangle
%     \State \textbf{Returns:} the handler responsible for processing the given command
%     \State \textbf{Description:} Determines the appropriate handler for the command based on its payload or type
% \EndProcedure
% \end{algorithmic}
% \end{algorithm}

% \subsection{CommandRouter}

% The \texttt{CommandRouter} interface defines the routing mechanism that maps each command to its corresponding handler.  
% It is responsible for ensuring that a command of type $\tau_P$ is dispatched only to a handler capable of processing that payload type $\tau_H$, following the eligibility constraint:

% \[
% \tau_H \sqsubseteq \tau_P \implies \text{handler is eligible for payload type } \tau_P
% \]

% Formally, the router defines a partial function:
% \[
% \rho : \mathcal{C}_{\tau_P} \rightarrow \mathcal{H}_{\tau_H}
% \]
% where $\mathcal{C}_{\tau_P}$ is the set of all command instances and $\mathcal{H}_{\tau_H}$ is the set of valid handlers.  
% If no suitable handler exists, $\rho$ may yield $\bot$ (undefined), triggering a \texttt{CommandNotRoutableException} or equivalent recovery mechanism.

% \begin{algorithm}[H]
% \caption{CommandRouter (\emph{Interface})}
% \begin{algorithmic}[1]
% \Procedure{Route}{command : Command\textlangle REQ\textrangle}
%     \State \textbf{Input:} command of type \textlangle REQ\textrangle
%     \State \textbf{Output:} CommandHandler of type \textlangle REQ, RES\textrangle
%     \State \textbf{Returns:} the handler responsible for processing the given command
%     \State \textbf{Description:} 
%         Determines the appropriate handler for the command based on its payload or type.
%         The routing process ensures type safety and deterministic dispatch:
%         \begin{itemize}
%             \item If multiple handlers are eligible, a priority or specificity rule is applied.
%             \item If no handler is found, a fallback or dead-letter route is invoked.
%             \item The resolution may be cached for repeated command types.
%         \end{itemize}
% \EndProcedure
% \end{algorithmic}
% \end{algorithm}

\subsection{CommandRouter}

The \texttt{CommandRouter} interface defines the routing mechanism that maps each command to a corresponding handler.  
It ensures that a command of type $\tau_P$ is dispatched only to a handler capable of processing that payload type $\tau_H$, following the eligibility constraint:

\[
\tau_H \sqsubseteq \tau_P \implies \text{handler is eligible for payload type } \tau_P
\]

Formally, the router defines a partial function:

\[
\rho : \mathcal{C}_{\tau_P} \rightarrow \mathcal{H}_{\tau_H}
\]

where $\mathcal{C}_{\tau_P}$ is the set of all command instances and $\mathcal{H}_{\tau_H}$ is the set of valid handlers.  
If no suitable handler exists, $\rho$ may yield $\bot$ (undefined), triggering a \texttt{CommandNotRoutableException} or an equivalent recovery mechanism.

\begin{algorithm}[H]
\caption{CommandRouter (\emph{Interface})}
\begin{algorithmic}[1]
\Procedure{Route}{command : Command\textlangle REQ\textrangle}
    \State \textbf{Input:} command of type \textlangle REQ\textrangle
    \State \textbf{Output:} CommandHandler of type \textlangle REQ, RES\textrangle
    \State \textbf{Returns:} the first matching handler responsible for processing the command
    \State \textbf{Description:} 
        Determines the appropriate handler based on the types of the request and expected response objects.
        \begin{itemize}
            \item The router scans the list of registered handlers and selects the first handler whose type matches the command payload.
            \item If multiple handlers are eligible, only the first match is returned.
            \item If no handler matches, a runtime exception or fallback mechanism is triggered.
        \end{itemize}
\EndProcedure
\end{algorithmic}
\end{algorithm}